\documentclass[12pt,a4paper,psamsfonts]{amsart}
\usepackage{amssymb,mathtools,calc}
\setlength{\topmargin}{0cm}
\setlength{\oddsidemargin}{0cm}
\setlength{\evensidemargin}{0cm}
\setlength{\textheight}{\paperheight - 2in - 35pt}
\setlength{\textwidth}{\paperwidth - 2in}
\setlength{\headheight}{12.5pt}
\setlength{\headsep}{25pt}
\setlength{\footskip}{30pt}
\renewcommand{\baselinestretch}{1.3}
\pagestyle{headings}
\raggedbottom
\setlength{\emergencystretch}{2em}
\usepackage[colorlinks=true,linkcolor=blue,urlcolor=blue]{hyperref}
\newtheorem{lemma}{Lemma}
\newtheorem{corollary}[lemma]{Corollary}
\DeclareMathOperator{\mult}{mult}
\DeclareMathOperator{\codim}{codim}
\title{Inequality on nef base locus}
\date{}
\begin{document}
\maketitle

For an effective Cartier divisor $F$ on a smooth variety and an
integral subvariety $V$, write $\mult_V(F)$ for the multiplicity
of $F$ at a general point of $V$. Equivalently, if $f$ is a local
equation of $F$ at the generic point $\eta_V$, then
\[
 \mult_V(F)=\max\{j\ge0:f\in\mathfrak m_{\eta_V}^{,j}\},
\]
where $\mathfrak m_{\eta_V}$ is the maximal ideal of
$\mathcal O_{X,\eta_V}$.

\begin{lemma}\label{lem:base-degree}
Let $X$ be a smooth projective variety of dimension $n$, and let
$D$ be a nef Cartier divisor. Let $\Lambda\subseteq|D|$ be a
nonempty linear system and $V$ an integral component of its base
locus, of codimension $c$. If, for an integer $h>0$,
\[
 \mult_V(F)\ge h\qquad\text{for every }F\in\Lambda,
\]
then
\[
 h^cD^{n-c}\cdot V\le D^n.
\]
\end{lemma}
\begin{proof}
Intersect successively with $c$ general members of $\Lambda$,
discarding the components contained in the base locus before
each further cut. The remaining components meet the next member
properly. Since $D$ is nef, discarding effective components does
not increase the degree measured by the remaining powers of $D$.
The resulting codimension-$c$ cycle therefore has $D$-degree at
most $D^n$.

At $\eta_V$, the base ideal is primary for the maximal ideal:
this uses that $V$ is a component of the base locus. No discarded
component contains $\eta_V$. The $c$ general local equations
$f_1,\ldots,f_c$ form a parameter ideal $I$ in the regular local
ring $R=\mathcal O_{X,\eta_V}$ of dimension $c$. Its local
intersection multiplicity is $\operatorname{length}(R/I)$.
The hypothesis gives $I\subseteq\mathfrak m^h$, so monotonicity
of Hilbert--Samuel multiplicity yields
\[
 \operatorname{length}(R/I)=e(I,R)
 \ge e(\mathfrak m^h,R)=h^c.
\]
Thus the coefficient of $V$ is at least $h^c$. Taking $D$-degree
and using nefness for all remaining effective components proves
the inequality.
\end{proof}

\begin{corollary}\label{cor:isolated-point}
Let $X,D,\Lambda$ be as in Lemma~\ref{lem:base-degree}, and let
$x$ be an isolated point of the base locus of $\Lambda$.
If $\mult_x(F)\ge r$ for every $F\in\Lambda$, where $r\ge0$, then
\[
 r^n\le D^n.
\]
\end{corollary}
\begin{proof}
For $r>0$, apply Lemma~\ref{lem:base-degree} with
$V=\{x\}$, $c=n$, and $h=\lceil r\rceil$.
For $r=0$, the assertion follows from nefness of $D$.
\end{proof}

\end{document}
